% GENERATED FILE. Edit canonical lessons, then run npm run generate:books.
% canonical-source-hash: fnv1a64:2526f385031f8bfa
% canonical-lessons: ML-C348-asrayiccirikkunnu, ML-C348-aghosham, ML-W348-gha, ML-C348-undayirunnenkil, ML-W348-six-lines, ML-R348-depends-and-would

\chapter{It Depends, and What Would Have Been}
\label{ch:ml-it-depends-and-what-would-have-been}
\hlchaptermodality{348}

\begin{chapteropening}
\textbf{By the end of this chapter:} \emph{I can say that something depends on something else, say what would happen or would have happened if things were different, and write six connected lines about a trip.}

Complete ML-R348-depends-and-would: it depends with āśrayiccirikkunnu, the unreal condition -irunneṅkil … -umāyirunnu, a celebration, and the new letter \ml{ഘ}.
\end{chapteropening}

\section[\texorpdfstring{\ml{അത്} \ml{കാലാവസ്ഥയെ} … (atŭ kālāvasthaye …)}{അത് കാലാവസ്ഥയെ … (atŭ kālāvasthaye …)}]{\ml{അത്} \ml{കാലാവസ്ഥയെ} \ml{ആശ്രയിച്ചിരിക്കുന്നു} (atŭ kālāvasthaye āśrayiccirikkunnu) — it depends on the weather}

\label{lesson:ML-C348-asrayiccirikkunnu}



\begin{quote}
Say \textquotedblleft{}the weather\textquotedblright{} and \textquotedblleft{}according to\textquotedblright{}. (\emph{kālāvastha}, \emph{anusariccŭ}.) Then \textquotedblleft{}if it rains, I will not come\textquotedblright{}. (\emph{maḻa peytāl ñān varilla.})
\end{quote}

\begin{grammarlens}[title={Grammar Lens: \ml{ആശ്രയിച്ചിരിക്കുന്നു}}]
What something depends on takes the object ending \textbf{-\ml{എ}}, as in \textbf{\ml{അവനെ}}, and the verb is \textbf{\ml{ആശ്രയിച്ചിരിക്കുന്നു}}: \textquotedblleft{}rests on, depends on\textquotedblright{}.

\begin{quote}
\textbf{\ml{നിങ്ങൾ} \ml{നാളെ} \ml{വരുമോ}? — \ml{അത്} \ml{കാലാവസ്ഥയെ} \ml{ആശ്രയിച്ചിരിക്കുന്നു}.}
\end{quote}

\begin{quote}
\emph{niṅṅaḷ nāḷe varumō? — atŭ kālāvasthaye āśrayiccirikkunnu.} — \textquotedblleft{}Will you come tomorrow? — It depends on the weather.\textquotedblright{}
\end{quote}

\textbf{\ml{സമയം}} gives \textbf{\ml{സമയത്തെ}}, and \textbf{\ml{പണം}} gives \textbf{\ml{പണത്തെ}}: \textquotedblleft{}it depends on time\textquotedblright{}, \textquotedblleft{}on money\textquotedblright{}. The verb is \textbf{\ml{ആശ്രയിക്കുക}}, \textquotedblleft{}to rely on\textquotedblright{}, built on Sanskrit \emph{āśraya}, \textquotedblleft{}that on which anything rests\textquotedblright{} (Monier-Williams), and its form here is the \textbf{-\ml{ഇച്ച്}} form, as in \textbf{\ml{അനുസരിച്ച്}}, with \textbf{\ml{ഇരിക്കുന്നു}} after it: it stays in a state of resting on.
\end{grammarlens}

\begin{grammarlens}[title={Grammar Lens: then both outcomes}]
After \textquotedblleft{}it depends\textquotedblright{}, give both sides:

\begin{quote}
\textbf{\ml{മഴ} \ml{പെയ്തില്ലെങ്കിൽ} \ml{ഞാൻ} \ml{വരും}. \ml{മഴ} \ml{പെയ്താൽ} \ml{വരില്ല}.}
\end{quote}

\begin{quote}
\emph{maḻa peytilleṅkil ñān varuṁ. maḻa peytāl varilla.} — \textquotedblleft{}If it does not rain, I will come. If it rains, I will not.\textquotedblright{}
\end{quote}

Everyday speech often reaches for the \textbf{\ml{അനുസരിച്ച്}} you know: \textbf{\ml{സമയം} \ml{അനുസരിച്ച്}}, \textquotedblleft{}depending on the time\textquotedblright{}.
\end{grammarlens}

\subsection*{Your turn}
Say these aloud:

\begin{itemize}
\raggedright
  \item \textquotedblleft{}It depends on the weather.\textquotedblright{}
  \item \textquotedblleft{}It depends on time\textquotedblright{}, then \textquotedblleft{}on money\textquotedblright{}
  \item both outcomes: if it does not rain, if it rains
  \item \textquotedblleft{}depending on the time\textquotedblright{} — \emph{samayaṁ anusariccŭ}
\end{itemize}

\begin{tcolorbox}[breakable,colback=teal!4,colframe=teal!35!black,title={Before you move on}]
\textquotedblleft{}It depends on the weather\textquotedblright{}? (\emph{\textbf{atŭ kālāvasthaye āśrayiccirikkunnu.}}) Which ending does the thing it depends on take? (\textbf{-\ml{എ}}: \emph{samayatte}, \emph{paṇatte}.)
\end{tcolorbox}
\section[\texorpdfstring{\ml{ആഘോഷം} (āghōṣaṁ)}{ആഘോഷം (āghōṣaṁ)}]{\ml{ആഘോഷം} (āghōṣaṁ) — a celebration}

\label{lesson:ML-C348-aghosham}



\begin{quote}
Before the new one: say \textquotedblleft{}Onam\textquotedblright{}, \textquotedblleft{}a festival\textquotedblright{}, and \textquotedblleft{}a wedding\textquotedblright{}. (\emph{ōṇaṁ}, \emph{utsavaṁ}, \emph{kalyāṇaṁ}.)
\end{quote}

\subsection*{You'll want to know: \ml{ആഘോഷം}}
\textbf{\ml{ആഘോഷം}} — \emph{āghōṣaṁ} — \textquotedblleft{}a celebration, festivities\textquotedblright{}.

\begin{quote}
\textbf{\ml{ഓണം} \ml{കേരളത്തിലെ} \ml{ഏറ്റവും} \ml{വലിയ} \ml{ആഘോഷമാണ്}.}
\end{quote}

\begin{quote}
\emph{ōṇaṁ kēraḷattile ēṟṟavuṁ valiya āghōṣamāṇŭ.} — \textquotedblleft{}Onam is the biggest celebration in Kerala.\textquotedblright{}
\end{quote}

\begin{quote}
\textbf{\ml{കല്യാണത്തിന്} \ml{വലിയ} \ml{ആഘോഷം} \ml{ഉണ്ടായിരുന്നു}.}
\end{quote}

\begin{quote}
\emph{kalyāṇattinŭ valiya āghōṣaṁ uṇṭāyirunnu.} — \textquotedblleft{}There was a big celebration at the wedding.\textquotedblright{}
\end{quote}

\textbf{\ml{ഉത്സവം}} is the festival itself, often at a temple; \textbf{\ml{ആഘോഷം}} is the celebrating, at a festival, a wedding or a birthday. The verb is \textbf{\ml{ആഘോഷിക്കുക}}, \textquotedblleft{}to celebrate\textquotedblright{}: \textbf{\ml{ഞങ്ങൾ} \ml{ഓണം} \ml{ആഘോഷിച്ചു}}.

Its first consonant is new. The next lesson puts it in your hand.

\begin{cousinweb}
\textbf{\ml{ആഘോഷം}} (\emph{āghōṣaṁ}) is a Sanskrit word, \emph{āghoṣa}, which Monier-Williams glosses as \textquotedblleft{}a calling out, a proclaiming\textquotedblright{}. In Malayalam it has come to mean the noise and gathering of a celebration; how the sense moved is not recorded there.
\end{cousinweb}

\subsection*{Your turn}
Say these aloud:

\begin{itemize}
\raggedright
  \item \emph{āghōṣaṁ}, with a breath after the \emph{g}
  \item \textquotedblleft{}Onam is the biggest celebration in Kerala.\textquotedblright{}
  \item \textquotedblleft{}There was a big celebration at the wedding.\textquotedblright{}
  \item \textquotedblleft{}We celebrated Onam.\textquotedblright{}
  \item \textquotedblleft{}Will you come to the celebration? — It depends on the time.\textquotedblright{} — \emph{āghōṣattinŭ varumō? — atŭ samayatte āśrayiccirikkunnu.}
\end{itemize}

\begin{tcolorbox}[breakable,colback=teal!4,colframe=teal!35!black,title={Before you move on}]
What does \textbf{\ml{ആഘോഷം}} mean? (\textbf{a celebration.}) How is it different from \textbf{\ml{ഉത്സവം}}? (\textbf{\ml{ഉത്സവം} is the festival; \ml{ആഘോഷം} is the celebrating.})
\end{tcolorbox}
\section[\texorpdfstring{\ml{ഘ} (gha)}{ഘ (gha)}]{\ml{ഘ} (gha) — the letter inside \ml{ആഘോഷം}}

\label{lesson:ML-W348-gha}



\begin{quote}
Before the new one: write \textbf{\ml{ഗ}} and \textbf{\ml{ഖ}} once each, and name them. (\emph{ga}, voiced; \emph{kha}, with a puff of breath.)

Now say \textbf{\ml{ആഘോഷം}} (\emph{āghōṣaṁ}), \textquotedblleft{}a celebration\textquotedblright{}. This lesson takes its first consonant out and puts it in your hand.
\end{quote}

\subsection*{Script you'll notice: \ml{ഘ}}
\textbf{\ml{ഘ}} — \emph{gha}: a \emph{g} with a puff of breath after it. It is to \textbf{\ml{ഗ}} what \textbf{\ml{ഖ}} is to \textbf{\ml{ക}}.

In \textbf{\ml{ആഘോഷം}} it wears the \emph{ō} sign, \textbf{\ml{ഘോ}}. Most words with \textbf{\ml{ഘ}} come from Sanskrit.

\subsection*{Writing: \ml{ഘ}}
Put your pen on \textbf{\ml{ഘ}} and follow its line. Copy the shape you can see, slowly, and larger than it is printed.

\begin{quote}
This book does not yet tell you \textbf{where to start this letter or which way to travel}. It is not written down here until it can be written down with a source. Copying what is in front of you needs no such source, and it is how the shape gets into your hand in the meantime.
\end{quote}

\subsection*{Your turn}
\begin{itemize}
\raggedright
  \item \emph{Find:} \ml{ഘ} inside \textbf{\ml{ആഘോഷം}}
  \item \emph{Trace it:} \ml{ഘ} three times, saying \emph{gh} as you finish each one
  \item \emph{Write it:} \ml{ഗ}, then \ml{ഘ}, from memory
  \item \emph{Write it:} \textbf{\ml{ആഘോഷം}}
\end{itemize}

\begin{tcolorbox}[breakable,colback=teal!4,colframe=teal!35!black,title={Before you move on}]
Which letter is this — \ml{ഘ}? (\textbf{gha}, \emph{g} with breath.) Which word did it come from? (\textbf{\ml{ആഘോഷം}}, \emph{āghōṣaṁ}, \textquotedblleft{}a celebration\textquotedblright{}.) Which letter is it the breathy partner of? (\textbf{\ml{ഗ}}.)
\end{tcolorbox}
\section[\texorpdfstring{\ml{പണം} … (paṇaṁ …)}{പണം … (paṇaṁ …)}]{\ml{ഉണ്ടായിരുന്നെങ്കിൽ} (uṇṭāyirunneṅkil) — if there had been: the unreal condition}

\label{lesson:ML-C348-undayirunnenkil}



\begin{quote}
Say \textquotedblleft{}I used to go\textquotedblright{}. (\emph{ñān pōkumāyirunnu}.) Then \textquotedblleft{}there was\textquotedblright{}. (\emph{uṇṭāyirunnu}.) Then \textquotedblleft{}if there is\textquotedblright{}. (\emph{uṇṭeṅkil}.)
\end{quote}

\begin{grammarlens}[title={Grammar Lens: -\ml{ഇരുന്നെങ്കിൽ}, if only it had been so}]
For something that is not so, or was not, put \textbf{\ml{ഇരുന്നെങ്കിൽ}} after the past form: \textbf{\ml{വന്നു}} $\to$ \textbf{\ml{വന്നിരുന്നെങ്കിൽ}}, \textquotedblleft{}if (he) had come\textquotedblright{}; \textbf{\ml{അറിഞ്ഞു}} $\to$ \textbf{\ml{അറിഞ്ഞിരുന്നെങ്കിൽ}}, \textquotedblleft{}if (I) had known\textquotedblright{}; \textbf{\ml{ഉണ്ടായി}} $\to$ \textbf{\ml{ഉണ്ടായിരുന്നെങ്കിൽ}}, \textquotedblleft{}if there had been\textquotedblright{}, \textquotedblleft{}if (I) had\textquotedblright{}.
\end{grammarlens}

\begin{grammarlens}[title={Grammar Lens: -\ml{ഉമായിരുന്നു}, would}]
The other half is the \textbf{-\ml{ഉമായിരുന്നു}} from the start of this core, now with its second job: \textquotedblleft{}would\textquotedblright{}, \textquotedblleft{}would have\textquotedblright{}.

\begin{quote}
\textbf{\ml{പണം} \ml{ഉണ്ടായിരുന്നെങ്കിൽ} \ml{ഞാൻ} \ml{പോകുമായിരുന്നു}.}
\end{quote}

\begin{quote}
\emph{paṇaṁ uṇṭāyirunneṅkil ñān pōkumāyirunnu.} — \textquotedblleft{}If I had money, I would go.\textquotedblright{}
\end{quote}

The same sentence also says \textquotedblleft{}if I had had money, I would have gone\textquotedblright{}; the context decides which.

\begin{quote}
\textbf{\ml{നീ} \ml{നേരത്തെ} \ml{വന്നിരുന്നെങ്കിൽ} \ml{നമുക്ക്} \ml{ബസ്} \ml{കിട്ടുമായിരുന്നു}.}
\end{quote}

\begin{quote}
\emph{nī nēratte vannirunneṅkil namukkŭ basŭ kiṭṭumāyirunnu.} — \textquotedblleft{}If you had come earlier, we would have caught the bus.\textquotedblright{}
\end{quote}

\begin{quote}
\textbf{\ml{അറിഞ്ഞിരുന്നെങ്കിൽ} \ml{ഞാൻ} \ml{വിളിക്കുമായിരുന്നു}.}
\end{quote}

\begin{quote}
\emph{aṟiññirunneṅkil ñān viḷikkumāyirunnu.} — \textquotedblleft{}If I had known, I would have called.\textquotedblright{}
\end{quote}

With \textbf{-\ml{അണം}} it is regret: \textbf{\ml{ഞാൻ} \ml{പോകണമായിരുന്നു}} (\emph{ñān pōkaṇamāyirunnu}), \textquotedblleft{}I should have gone\textquotedblright{}.
\end{grammarlens}

\subsection*{Your turn}
Say these aloud:

\begin{itemize}
\raggedright
  \item \textquotedblleft{}If I had money, I would go.\textquotedblright{}
  \item \textquotedblleft{}If you had come earlier, we would have caught the bus.\textquotedblright{}
  \item \textquotedblleft{}If I had known, I would have called.\textquotedblright{}
  \item \textquotedblleft{}I should have gone.\textquotedblright{}
\end{itemize}

\begin{tcolorbox}[breakable,colback=teal!4,colframe=teal!35!black,title={Before you move on}]
\textquotedblleft{}If he had come\textquotedblright{}? (\emph{vannirunneṅkil}.) \textquotedblleft{}Would have called\textquotedblright{}? (\emph{\textbf{viḷikkumāyirunnu.}}) How does \emph{pōkumāyirunnu} tell \textquotedblleft{}used to go\textquotedblright{} from \textquotedblleft{}would have gone\textquotedblright{}? (\textbf{By the context: an if-clause in front makes it \textquotedblleft{}would\textquotedblright{}.})
\end{tcolorbox}
\section[(six lines)]{Six lines about a trip}

\label{lesson:ML-W348-six-lines}



\begin{quote}
The last A2 writing task asked for a connected message about a familiar day. This one asks for six lines, and they tell a story.

Pick one reader: a friend, by name.
\end{quote}

\subsection*{What to know first}
A told story is not a list of events. It has a \textbf{scene}, in \textbf{\ml{ആയിരുന്നു}} and the \textbf{-\ml{ഉകയായിരുന്നു}} past. It has \textbf{events} in the plain past, kept in order with \textbf{\ml{ആദ്യം}}, \textbf{\ml{എന്നിട്ട്}}, \textbf{-\ml{അപ്പോൾ}}, \textbf{\ml{അതിനുശേഷം}} and \textbf{\ml{ഒടുവിൽ}}. It has at least one \textbf{reason}, with \textbf{-\ml{അതുകൊണ്ട്}} or \textbf{\ml{കാരണം}}, and a \textbf{feeling}, with \textbf{\ml{എനിക്ക്} … \ml{തോന്നി}}.

A good one ends by looking at what might have been: \textbf{… -\ml{ഇരുന്നെങ്കിൽ} … -\ml{ഉമായിരുന്നു}}.

\subsection*{Writing — connected composition}
On paper, write \textbf{six lines} to your reader about a trip, real or invented, one line for each of these:

\begin{enumerate}
\raggedright
  \item the scene: when it was, with \textbf{\ml{ഈയിടെ}} or \textbf{\ml{കഴിഞ്ഞ}}, and what it was like, with \textbf{\ml{ആയിരുന്നു}};
  \item where you went and what you were doing on the way, in the \textbf{-\ml{ഉകയായിരുന്നു}} past;
  \item the first thing that happened, with \textbf{\ml{ആദ്യം}};
  \item what happened next, with \textbf{\ml{എന്നിട്ട്}} or \textbf{-\ml{അപ്പോൾ}}, and one thing that went wrong;
  \item how it ended and how you felt, with \textbf{\ml{ഒടുവിൽ}} and \textbf{\ml{എനിക്ക്} … \ml{തോന്നി}}, and a reason;
  \item one line about what might have been, with \textbf{-\ml{ഇരുന്നെങ്കിൽ} … -\ml{ഉമായിരുന്നു}}.
\end{enumerate}

There is no model on this page. Use only words you already write. The whole exercise is choosing the right past for each line, and joining them.

\begin{tcolorbox}[breakable,colback=teal!4,colframe=teal!35!black,title={Before you move on}]
Now read your lines as you read the village story earlier in this core. Mark every verb \textbf{s} for scene or \textbf{e} for event. Line two should be mostly \textbf{s}; lines three and four mostly \textbf{e}.

Then three spelling checks. A word that starts with long \emph{ī} starts with \textbf{\ml{ഈ}}, not \textbf{\ml{ഇ}}. The \emph{au} of \textbf{\ml{സൗകര്യം}} is the tail \textbf{\ml{ൗ}} after its consonant. \textbf{\ml{ദുഃഖം}} keeps its two rings, \textbf{\ml{ഃ}}.

Last, read line six alone. Does the if-clause end in \textbf{-\ml{ഇരുന്നെങ്കിൽ}}, and the main verb in \textbf{-\ml{ഉമായിരുന്നു}}?
\end{tcolorbox}
\section[(review)]{Review — it depends, and what would have been}

\label{lesson:ML-R348-depends-and-would}



\begin{quote}
Say \textquotedblleft{}a celebration\textquotedblright{}. (\emph{āghōṣaṁ}.) Write \textbf{\ml{ഘ}} once from memory, then \textbf{\ml{ആഘോഷം}}.
\end{quote}

\subsection*{Your turn}
\begin{itemize}
\raggedright
  \item \emph{Say it:} \textquotedblleft{}Will you come tomorrow? — It depends on the weather.\textquotedblright{}
  \item \emph{Say it:} \textquotedblleft{}If it does not rain, I will come. If it rains, I will not.\textquotedblright{}
  \item \emph{Say it:} \textquotedblleft{}Onam is the biggest celebration in Kerala.\textquotedblright{}
  \item \emph{Say it:} \textquotedblleft{}If I had money, I would go.\textquotedblright{}
  \item \emph{Say it:} \textquotedblleft{}If you had come earlier, we would have caught the bus.\textquotedblright{}
  \item \emph{Say it:} \textquotedblleft{}If I had known, I would have called. I should have gone.\textquotedblright{}
  \item \emph{Write it:} \textbf{\ml{ആഘോഷം}}
\end{itemize}

\begin{tcolorbox}[breakable,colback=teal!4,colframe=teal!35!black,title={Before you move on}]
\textquotedblleft{}It depends on …\textquotedblright{}? (\textbf{… -\ml{എ} \ml{ആശ്രയിച്ചിരിക്കുന്നു}}.) \textquotedblleft{}If I had known\textquotedblright{}? (\textbf{\ml{അറിഞ്ഞിരുന്നെങ്കിൽ}}.) \textquotedblleft{}Would have called\textquotedblright{}? (\textbf{\ml{വിളിക്കുമായിരുന്നു}}.) \textquotedblleft{}A celebration\textquotedblright{}? (\textbf{\ml{ആഘോഷം}}, with \textbf{\ml{ഘ}}.)
\end{tcolorbox}
